% Ad Atticum 14.19 --- headnote
% ad-atticum-14-19, 15 May 44 BC, in Pompeiano

\textit{Cicero to Atticus, written at the Pompeian villa on 15 May
44 BC --- Perseus dateline \textit{Scr. in Pompeiano i Id. Mai. a.
710 (44)} (the day before the Ides). Two of Atticus's letters
have caught up with Cicero at once. He works through them in
order: relief that his own letter to Cassius crossed Atticus's
recommending the same line, despair at Dolabella's unpaid debt
(the manuscripts carry the corruption \textit{†aritia†} here),
and then surprise at finding two consoling letters, one from
Brutus himself and one from Atticus, telling him Brutus is
meditating exile. Cicero's harbour, he says, is a different
one --- one nearer for a man his age. He and Atticus agree:
this is no season for a man of their years to take to a camp,
\textit{praesertim civilibus}.}

\textit{The second half is a tour of the political weather.
Antony is gracious about the Clodius affair, allegedly furious
with Pansa about Clodius and Deiotarus, and --- unattractively
--- disapproves of Dolabella's clearing of the Forum column.
The young man who wore a garland for Caesar and put it off in
mourning was Atticus's nephew Quintus, defending himself to his
father. Cicero has written carefully to Dolabella (in the
spirit Atticus prescribed) and to Sicca; he holds the line with
Servius's caution; he closes the door politely on Publilius's
overtures via Caerellia, on the grounds that what is being
asked is not even open to him. Greek appears twice in section
5, both as \textit{[Greek: praxin]} (``exploit''), and once at
the end as the agreed gauge of Dolabella's performance: the
greater \textit{praxis} will be settling the debt.}
